\section*{An explicit all-parameter rare-bias tail for the two-bit half-tax inequality}

All logarithms in the proof are natural.  Let
\[
 0<p\le {1\over64},\qquad 0\le q\le1,
\]
and put
\[
 s=4q(1-q),\quad c=4q^2,\quad b=1-sp,\quad
 a={cp(1-p)\over b}.
\]
For $h(x)=-x\log x-(1-x)\log(1-x)$, define
\[
 d(p,q)=b h(a)+sp h(p)-h(sp)+p h(s)
          +{p\over2}\log{1-p\over p}.                 \tag{1}
\]
Thus $d=(\log 2)D$, where
\[
 D=M_A+\tfrac12\Delta_A-J_A
\]
is the bit-valued margin in the frozen two-bit AND candidate.

\paragraph{Theorem.}
For every $0<p\le1/64$ and every $q\in[0,1]$,
\[
 \boxed{d(p,q)>0.}                                     \tag{2}
\]
Equivalently, $J_A<M_A+\Delta_A/2$ throughout the complete rare-bias
tail, uniformly in the second bias.

\paragraph{Removing the singular parameter.}
Write
\[
 \ell(x)=-(1-x)\log(1-x),\qquad
 \psi(x)={\ell(x)\over x}\ (x>0),\quad\psi(0)=1,
\]
and $\phi(x)=x\log(1/x)$, with $\phi(0)=0$.  Expanding
$h(x)=x\log(1/x)+\ell(x)$ in (1), using $ba=cp(1-p)$, and cancelling the
$s\log s$ terms gives the exact identity
\[
 {d(p,q)\over p}=A(p,q)\log{1\over p}+F(p,q),           \tag{3}
\]
where
\[
 A(p,q)=8(q-\tfrac14)^2+p(s-c),                         \tag{4}
\]
and
\begin{align*}
 F(p,q)={}&c(1-p)\log{b\over1-p}+(1-p)\phi(c)
       +c(1-p)\psi(a)\\
 &+sp\psi(p)-s\psi(sp)+\ell(s)+\tfrac12\log(1-p).
                                                               \tag{5}
\end{align*}
Every term in (5) extends continuously to $p=0$; in particular there is
no numerical $0/0$ endpoint.

The coefficient (4) is nonnegative on the full closed rectangle.  For
$q\le1/2$, both terms in (4) are nonnegative.  For $q\ge1/2$, the second
term is nonpositive, so its worst value occurs at $p=1/64$, and
\[
 16A(p,q)\ge128(q-\tfrac14)^2+q(1-2q)
             =126q^2-63q+8>0,                          \tag{6}
\]
because the last quadratic has discriminant $63^2-4(126)(8)=-63$.
Since $\log(1/p)\ge\log64$, (3)--(6) imply
\[
 {d(p,q)\over p}\ge E_{64}(p,q):=A(p,q)\log64+F(p,q). \tag{7}
\]
The right side is continuous on the compact rectangle
$[0,1/64]\times[0,1]$.  This is the analytic tail/compactness reduction;
the interval computation below certifies this continuous remainder, not
the singular original expression.

\paragraph{Directed-rounding certificate.}
The accompanying program encloses $E_{64}$ using 80-decimal Decimal
arithmetic with ROUND\_FLOOR for every lower endpoint and ROUND\_CEILING
for every upper endpoint.  Besides interval arithmetic, it uses only the
following proved one-variable facts:
\[
 \psi'(x)={\log(1-x)+x\over x^2}\le0,
\]
$\ell$ has its unique maximum $1/e$ at $1-1/e$, and $\phi$ has its unique
maximum $1/e$ at $1/e$.  The functions $s=4q(1-q)$ and $c=4q^2$ are
enclosed with their exact monotonicity.  The known inequality $A\ge0$ from
(6) is used only to intersect its interval enclosure with $[0,\infty)$.

Adaptive bisection visited 175 boxes, accepted 88 leaves, reached depth 8,
and left no unresolved box.  Its smallest certified lower endpoint is
\[
 E_{64}>0.0116100600538253879943401888136485579611459563.
\]
Together with (7), this proves (2).  No floating-point value is used in a
sign decision.

\paragraph{Exact scope and GU relevance.}
This theorem removes the formerly nonexplicit joining radius in the
$p\downarrow0$ argument for the two-bit Boolean AND kernel.  It is not a
universal GU theorem.  In any proposed GU assembly where the local
restriction mismatch contains this kernel, it permits the sharp payment
$J_A-M_A\le\Delta_A/2$ uniformly down to deterministic bias.  A universal
conclusion still needs (i) the corresponding bound for every higher-arity
Boolean block and (ii) a proof that the $M_A$ terms cancel or telescope and
that the physical $\Delta_A/2$ taxes are not reused.  Under precisely those
two additional assembly hypotheses, summing the local inequalities gives
the displayed sufficient implication
\[
 \sum J_A\le\sum M_A+\tfrac12\sum\Delta_A,quad
 \sum M_A\text{ cancelled},\quad \sum\Delta_A\le I(f)
 \quad\Longrightarrow\quad G(f)\le\tfrac12 I(f).       \tag{8}
\]
Neither higher arity nor the two global single-use premises in (8) is
proved here; full FEI remains open.

